\section{Executable pebble game}
\label{sec:executable}

The previous chapter (\cref{sec:pebble-game}) formalises the basic
$(k, \ell)$-pebble game and proves its certificate-form correctness
theorem \cref{thm:pebble-game-correct}: in the matroidal regime
$\ell < 2k$, a finite simple graph $G$ is $(k, \ell)$-sparse if and
only if the pebble game accepts. But the convenience form
\texttt{runPebbleGame} is \texttt{noncomputable}
(\cref{fmlnote:pebble-game-computable-core}): its body enumerates the
finite sets it works over as lists through mathlib operations
(\texttt{Finset.toList} on a partial orientation's out-neighbours,
\texttt{Quot.out} on the edge set) that reduce only through
\texttt{Classical} choice, so it can be neither evaluated by
\texttt{\#eval} nor reduced by \texttt{native\_decide}.

This chapter closes the gap. The pebble game was already factored into
the computable core (\texttt{tryReachPebbleWith} /
\texttt{tryAddEdgeWith} / \texttt{runPebbleGameWith}) of
\cref{fmlnote:pebble-game-computable-core}; here the core is run
through a computable form \texttt{runPebbleGameExec} that replaces the
two noncomputable list enumerations with \texttt{Finset.sort}-based
analogues under a linear order on $V$. The certificate-form
correctness theorem lifts to \texttt{runPebbleGameExec} with no new
mathematical content; on top of it we register $\mathrm{Decidable}$
instances for the project predicates $\mathrm{IsSparse}\,k\,\ell$,
$\mathrm{IsTight}\,k\,\ell$, and $\mathrm{IsLaman}$. The instances
reduce in polynomial time; they decide the same propositions a
brute-force iteration over all finite subsets of $V$ would decide, but
choose a faster reduction.

A small \texttt{lake exe pebble-game} CLI built on the same
\texttt{Decidable} instance reads an edge-list file and prints
$\mathrm{LAMAN}$ / $\mathrm{SPARSE\_NOT\_TIGHT}$ /
$\mathrm{NOT\_SPARSE}$.

\subsection{Runtime and backends}
\label{sec:executable-backends}

The same function body can be run through three
syntactically-different invocations, each executed by a different
runtime. The three invocations that reduce
\texttt{runPebbleGameExec}'s compiled body are:

\begin{itemize}
  \item \texttt{\#eval (decide (G.IsSparse 2 3))} (top-level
    elaborator command). Reduces $\mathrm{decide}\,P$ through the
    registered $\mathrm{Decidable}$ instance, whose body is the
    pebble game; runs in the bytecode interpreter. Fast in practice
    (polynomial in $|V| + |E|$, no kernel small-step overhead).
  \item \texttt{lake exe pebble-game input.txt} (built CLI binary).
    Same compiled body, but executed by Lean's native-code backend
    rather than the bytecode interpreter. Strictly faster constant
    factor; identical big-$O$.
  \item \texttt{by decide} (tactic, inside a proof of the form
    \texttt{theorem foo : G.IsLaman := by decide}). Reduces the
    \emph{same} $\mathrm{Decidable}$ instance, but through the
    Lean kernel's small-step reducer, not the compiled backend.
    Polynomially many small-step reductions per pebble-game
    operation $\Rightarrow$ impractically slow on any
    realistic graph. The tactic
    \texttt{native\_decide} compiles and runs the same body
    natively, paying a trust-the-compiled-binary axiom (Lean's
    \texttt{nativeDecide}); its trust footprint is identical to
    the CLI's.
\end{itemize}

Only \texttt{\#eval} and \texttt{lake exe pebble-game} are polynomial
and end-to-end on real graphs; \texttt{by decide} stays formally
correct but its runtime is academic. The worked examples
(\cref{sec:executable-examples}) use \texttt{\#eval} or
\texttt{native\_decide} accordingly.

For \texttt{\#eval} and \texttt{native\_decide} to reduce, the
registered $\mathrm{Decidable}$ instance body must be free of
noncomputable primitives; this is why \texttt{runPebbleGameExec} is
built afresh rather than reused from \texttt{runPebbleGame}
(\cref{fmlnote:executable-computable-core}).

\subsection{Computable list views}
\label{sec:executable-list-views}

The two noncomputable list enumerations in \texttt{runPebbleGame} are
localised: an out-neighbour list of a partial orientation $D$ at a
vertex $v$ (the input to the verified depth-first search
\texttt{reachableFinding}), and an enumeration of the input graph's
edges as ordered pairs (the input to the \texttt{runPebbleGameWith}
fold). Both are made computable by $\mathrm{Finset.sort}$ under a
linear order on $V$.

\begin{definition}[Sorted out-neighbour list]
  \label{def:outListSorted}
  \lean{CombinatorialRigidity.PebbleGame.PartialOrientation.outListSorted}
  \leanok
  \uses{def:partial-orientation, def:pebble-counts}
  Let $V$ be a finite type with a linear order and decidable
  equality, and let $D$ be a partial orientation on $V$
  (\cref{def:partial-orientation}). For $v \in V$, define
  \[
    \mathrm{outListSorted}_D(v)
    \;:=\; (D.\mathrm{outNbhd}\,v).\mathrm{sort}\,(\cdot \le \cdot)
    \quad : \quad \mathrm{List}\,V,
  \]
  the out-neighbours of $v$ in $D$ listed in increasing order. It is
  the computable replacement for the noncomputable $\mathrm{outList}$
  of $D$.
\end{definition}

\begin{lemma}[Membership in the sorted out-neighbour list]
  \label{lem:mem-outListSorted}
  \lean{CombinatorialRigidity.PebbleGame.PartialOrientation.mem_outListSorted}
  \leanok
  \uses{def:outListSorted}
  Under the hypotheses of \cref{def:outListSorted}, for all $v, u
  \in V$,
  \[
    u \in \mathrm{outListSorted}_D(v) \;\iff\; (v, u) \in D.\mathrm{arcs}.
  \]
\end{lemma}
\begin{proof}
  \leanok
  $\mathrm{Finset.sort}$ preserves membership; the result follows
  from $D.\mathrm{mem\_outNbhd}$.
\end{proof}

\begin{definition}[Sorted edge enumeration]
  \label{def:edgeListSorted}
  \lean{SimpleGraph.edgeListSorted}
  \leanok
  \uses{def:partial-orientation}
  Let $V$ have a linear order and decidable equality, and let $G$ be a
  finite simple graph on $V$. Define
  \[
    \mathrm{edgeListSorted}\,G
    \;:=\; \mathrm{ofLex}\,\circ\,
      \big((G.\mathrm{edgeFinset}.\mathrm{image}\,(e \mapsto
        \mathrm{toLex}\,(e.\inf, e.\sup))).\mathrm{sort}\,(\cdot \le \cdot)\big)
    \quad : \quad \mathrm{List}\,(V \times V),
  \]
  where each edge $e$ is first imaged to its lex-smaller-first ordered
  pair $(e.\inf, e.\sup)$ viewed as $\mathrm{Lex}\,(V \times V)$, then
  sorted under $\mathrm{Prod.Lex}$'s linear order, and finally untagged
  via $\mathrm{ofLex}$ to recover $V \times V$. It is the computable
  replacement for the noncomputable enumeration of $E(G)$ used by
  \texttt{runPebbleGame}.
\end{definition}

\begin{fmlnote}
  Sorting the edge set $G.\mathrm{edgeFinset}$ directly by
  $(\cdot \le \cdot)$ is not available: mathlib already fixes the order
  on $\mathrm{Sym}_2 V$ to be the non-total subset order coming from
  its set-like structure, so there is no total order on
  $\mathrm{Sym}_2 V$ to sort against. Sorting the projection
  $(e.\inf, e.\sup)$ under the lexicographic order on $V \times V$ is
  the idiomatic alternative and needs no new order instance.
\end{fmlnote}

\begin{lemma}[Membership in the sorted edge list]
  \label{lem:mem-edgeListSorted}
  \lean{SimpleGraph.mem_edgeListSorted}
  \leanok
  \uses{def:edgeListSorted}
  Under the hypotheses of \cref{def:edgeListSorted}, for all
  $u, v \in V$,
  \[
    (u, v) \in \mathrm{edgeListSorted}\,G \;\iff\;
      u \le v \;\wedge\; \{u, v\} \in G.\mathrm{edgeFinset}.
  \]
\end{lemma}
\begin{proof}
  \leanok
  Forward: an entry $(u, v) \in \mathrm{edgeListSorted}\,G$ is
  $\mathrm{ofLex}(\mathrm{toLex}\,(e.\inf, e.\sup))$ for some
  $e = \{a, b\} \in G.\mathrm{edgeFinset}$, so $(u, v) = (e.\inf,
  e.\sup)$; splitting on whether $a \le b$ or $b \le a$ identifies
  $\{u, v\} = e$, whence $\{u, v\} \in G.\mathrm{edgeFinset}$.
  Backward: given $u \le v$ and $\{u, v\} \in G.\mathrm{edgeFinset}$,
  take $e := \{u, v\}$; then $e.\inf = u$ and $e.\sup = v$ since
  $u \le v$.
\end{proof}

\subsection{Computable form}
\label{sec:executable-wrapper}

The algorithm was already factored into the computable core
\texttt{runPebbleGameWith}, parametrised over a caller-supplied
out-neighbour enumeration and edge list; the convenience form
\texttt{runPebbleGame} supplies the noncomputable enumerations
$\mathrm{outList}$ and \texttt{G.edgeFinset.toList.map Quot.out}. We
now add a computable sibling that supplies \cref{def:outListSorted}
and \cref{def:edgeListSorted} instead.

\begin{definition}[Computable verdict-bearing pebble game]
  \label{def:runPebbleGameExec}
  \lean{CombinatorialRigidity.PebbleGame.PartialOrientation.runPebbleGameExec}
  \leanok
  \uses{def:runPebbleGame, def:outListSorted, def:edgeListSorted,
    def:pebbleGameResult, thm:pebble-game-correct,
    lem:workhorseWitness-certifies}
  Let $V$ be a finite type with a linear order and decidable equality,
  let $G$ be a finite simple graph on $V$, and let $\ell < 2k$. Define
  \[
    \mathrm{runPebbleGameExec}\,G\,k\,\ell : \mathrm{PebbleGameResult}\,G\,k\,\ell
  \]
  by running the computable core $\mathrm{runPebbleGameWith}$ from the
  empty orientation with the sorted list views \cref{def:outListSorted}
  and \cref{def:edgeListSorted}, then packaging its $\mathrm{Sum}$
  output as a $\mathrm{PebbleGameResult}$ verdict
  (\cref{def:pebbleGameResult}): an accepting orientation $D$ (the
  $\mathrm{.inr}$ branch) becomes $\mathrm{.accept}\,D$, and a failure
  witness $w$ (the $\mathrm{.inl}$ branch, \cref{def:workhorseWitness})
  becomes $\mathrm{.reject}\,w.V'$. It is computable end-to-end and has
  no $\mathrm{Classical}$ dependencies.
\end{definition}

\begin{fmlnote}
  \label{fmlnote:executable-computable-core}
  The convenience form $\mathrm{runPebbleGame}$ is noncomputable
  because it enumerates finite sets as lists through mathlib operations
  that reduce only through $\mathrm{Classical}$ choice
  (\cref{fmlnote:pebble-game-computable-core}): $\mathrm{Finset.toList}$
  on the out-neighbours and $\mathrm{Quot.out}$ on the edge set.
  Rebuilding the same algorithm on the $\mathrm{Finset.sort}$-based
  views \cref{def:outListSorted} and \cref{def:edgeListSorted} removes
  both, so $\mathrm{runPebbleGameExec}$ reduces with no
  $\mathrm{Classical}$ dependency --- the requirement for
  $\mathrm{\#eval}$ and $\mathrm{native\_decide}$ to run its registered
  $\mathrm{Decidable}$ instance.

  The verdict repackaging carries proofs but adds no computational
  content: on the accept branch the two obligations of
  \cref{def:pebbleGameResult} --- that the orientation's underlying
  edge set is $E(G)$ and that it is $(k, \ell)$-reachable --- are
  discharged by \cref{thm:pebble-game-correct}, and on the reject
  branch the size and strict-inequality obligations by
  \cref{lem:workhorseWitness-certifies}. The verdict is thus structural
  metadata over the core's raw output, and the compiled body that
  $\mathrm{\#eval}$ reduces is the pebble game itself.
\end{fmlnote}

The computable form inherits the correctness theorem
\cref{thm:pebble-game-correct} with no new mathematical content. That
theorem is parametrised over the out-neighbour enumeration and the
candidate edge list; the only step specific to \cref{def:edgeListSorted}
is checking that it meets the theorem's three side conditions on the
edge list --- loop-freeness, pairwise distinctness of the underlying
unordered edges, and that their $\mathrm{Sym}_2$-image is exactly
$G.\mathrm{edgeFinset}$ --- which the accompanying discharge lemmas
establish. Its verdict-form analog
\cref{thm:pebbleGameResult-isAccept-iff} restates it with the
certificate equivalence absorbed into the verdict's type.

\subsection{Decidability of sparsity, tightness, and Laman}
\label{sec:executable-decidable}

\begin{definition}[\texttt{Decidable} instance for sparsity]
  \label{def:isSparse-decidable}
  \lean{SimpleGraph.instDecidableIsSparse}
  \leanok
  \uses{def:runPebbleGameExec, def:isSparse,
    thm:pebbleGameResult-isAccept-iff}
  Let $V$ be a finite type with a linear order and decidable equality,
  and let $G$ be a finite simple graph on $V$. In the matroidal regime
  $\ell < 2k$ --- carried in the formalisation as a
  $\mathrm{Fact}\,(\ell < 2k)$ instance for typeclass synthesis --- the
  proposition $G.\mathrm{IsSparse}\,k\,\ell$ is decidable via the rule
  \[
    \mathrm{decide}\,(G.\mathrm{IsSparse}\,k\,\ell)
    \;=\;
    (\mathrm{runPebbleGameExec}\,G\,k\,\ell).\mathrm{isAccept},
  \]
  the accepting-verdict projection of the computable pebble game. It is
  the canonical $\mathrm{Decidable}$ instance for
  $\mathrm{IsSparse}\,k\,\ell$; because the verdict is structural
  metadata (\cref{fmlnote:executable-computable-core}), the compiled
  body it reduces is the pebble game itself.
\end{definition}

\begin{remark}[Polynomial vs.\ exponential reduction]
  \label{rem:decidable-runtime}
  $G.\mathrm{IsSparse}\,k\,\ell$ quantifies over all finite subsets of
  $V$, so it also admits a brute-force decision procedure that iterates
  over all $2^{|V|}$ subsets (via
  $\mathrm{Fintype.decidableForallFintype}$). Both procedures decide the
  \emph{same} proposition and produce a valid $\mathrm{Decidable}$
  instance; they differ only in the instance body, which is what
  $\mathrm{\#eval}$ and $\mathrm{decide}$ reduce. The pebble-game
  instance reduces in polynomial time, the brute-force instance in
  exponential time. To keep typeclass search from selecting the slow
  path, the project registers exactly one $\mathrm{Decidable}$ instance
  per predicate.
\end{remark}

\begin{definition}[\texttt{Decidable} instance for tightness]
  \label{def:isTight-decidable}
  \lean{SimpleGraph.instDecidableIsTight}
  \leanok
  \uses{def:isSparse-decidable, def:isTight}
  Under the hypotheses of \cref{def:isSparse-decidable},
  $G.\mathrm{IsTight}\,k\,\ell$ is decidable: it is the conjunction of
  $G.\mathrm{IsSparse}\,k\,\ell$ (\cref{def:isSparse-decidable}) with
  the edge-count equation $|E(G)| + \ell = k|V|$, and the latter is
  decidable as an equality of natural numbers. (The definition of
  $\mathrm{IsTight}$ states this equation with cardinalities as
  $\mathrm{Set.ncard}$ and $\mathrm{Nat.card}$; the formalisation
  rewrites it to the $\mathrm{Finset.card}$ and $\mathrm{Fintype.card}$
  form that reduces, a step with no algorithmic content.)
\end{definition}

\begin{definition}[\texttt{Decidable} instance for Laman]
  \label{def:isLaman-decidable}
  \lean{SimpleGraph.instDecidableIsLaman}
  \leanok
  \uses{def:isTight-decidable, def:isLaman}
  Under the hypotheses of \cref{def:isSparse-decidable} at
  $(k, \ell) = (2, 3)$, $G.\mathrm{IsLaman}$ is decidable: it is
  $G.\mathrm{IsTight}\,2\,3$ (\cref{def:isTight-decidable}) unfolded.
  The formalisation registers the matroidal-regime witness
  $\mathrm{Fact}\,(3 < 2 \cdot 2)$ as a top-level instance, so the
  Laman case is zero-hypothesis at the call site:
  $\mathrm{\#eval}\,(\mathrm{decide}\,G.\mathrm{IsLaman})$ runs without
  any caller-supplied hypothesis.
\end{definition}

\subsection{Worked examples}
\label{sec:executable-examples}

The end-to-end-executability target is exercised by a small set of
$\mathrm{Fin}\,n$ examples in
\texttt{CombinatorialRigidity/PebbleGame/Examples.lean}. Each is a
$\mathrm{\#eval}$ of $\mathrm{decide}\,G.\mathrm{IsLaman}$ on a
concrete graph; the output should match the well-known
classification.

\begin{enumerate}
  \item \textbf{$K_4 \setminus e$} on $\mathrm{Fin}\,4$. Laman: 4
    vertices, 5 edges, $2 \cdot 4 - 3 = 5$. The worked example of
    \cref{thm:top-fin-four-minus-edge-isLaman}; verifies
    the minimal non-trivial accept path.
  \item \textbf{Moser spindle.} A 7-vertex 11-edge Laman graph
    ($2 \cdot 7 - 3 = 11$). Verifies the accept path on a graph
    large enough for the depth-first search to run non-trivially.
  \item \textbf{$K_4$ on $\mathrm{Fin}\,4$.} Not Laman, globally
    over-constrained: 4 vertices, 6 edges, $6 > 2 \cdot 4 - 3 = 5$.
    The pebble game rejects on the full-vertex subset at the
    sparsity step; the blocking witness
    (\cref{def:blockingWitness}) $V' = V$ violates Invariant~(4).
  \item \textbf{A sparse-but-not-tight example.} A path on 5
    vertices: $(2, 3)$-sparse but $|E| = 4 < 2 \cdot 5 - 3 = 7$,
    so not tight, hence not Laman. The pebble game accepts; the
    \cref{def:isTight-decidable} edge-count check rejects.
\end{enumerate}

These examples additionally demonstrate the
$\mathrm{\#eval}$ / $\mathrm{native\_decide}$ runtime claim of
\cref{sec:executable-backends}: the reduction completes in
sub-second wall time on a developer laptop.

\subsection{CLI binary}
\label{sec:executable-cli}

A \texttt{lake exe pebble-game} entry point reads a simple
edge-list file and prints the classification together with the
inline witness data carried by the \cref{def:pebbleGameResult}
verdict. Input format: a leading line
\texttt{n m} giving vertex count and edge count, followed by $m$
lines of two whitespace-separated integers $0 \le u, v < n$. Output
schema:
\begin{center}
\begin{tabular}{lll}
$\mathrm{LAMAN}$           & $\mathrm{SPARSE\_NOT\_TIGHT}$ & $\mathrm{NOT\_SPARSE}$ \\
$\mathrm{ARCS}\ u_0\ v_0$  & $\mathrm{ARCS}\ u_0\ v_0$     & $\mathrm{BLOCKING}\ |V'|$ \\
$\mathrm{ARCS}\ u_1\ v_1$  & $\mathrm{ARCS}\ u_1\ v_1$     & $\mathrm{VERTEX}\ w_0$ \\
$\vdots$                   & $\vdots$                      & $\mathrm{VERTEX}\ w_1$ \\
                           &                               & $\vdots$ \\
\end{tabular}
\end{center}
i.e.\ a verdict label on the first line, then witness lines: on
$\mathrm{LAMAN}$ / $\mathrm{SPARSE\_NOT\_TIGHT}$, one
$\mathrm{ARCS}\ u\ v$ line per arc $(u, v) \in D.\mathrm{arcs}$ of
the accepting partial orientation, in lexicographic order on $(u, v)$;
on $\mathrm{NOT\_SPARSE}$, a $\mathrm{BLOCKING}\ |V'|$ line giving the
blocking subset's size followed by $|V'|$ lines $\mathrm{VERTEX}\ w$,
one per $w \in V'$, in increasing order.

The binary is a thin \texttt{IO} program:
\begin{enumerate}
  \item Parse the input file; reject malformed input with a
    diagnostic.
  \item Construct $G : \mathrm{SimpleGraph}\,(\mathrm{Fin}\,n)$
    from the parsed edge list via $\mathrm{fromEdgeSet}$.
  \item Invoke $\mathrm{runPebbleGameExec}\,G\,2\,3$ directly to
    obtain a \cref{def:pebbleGameResult} verdict (the
    \cref{def:isLaman-decidable} instance is no longer needed at the
    CLI's entry point --- the Boolean projection it derives via
    \cref{thm:pebbleGameResult-isAccept-iff} is recovered from the
    verdict's constructor name).
  \item Pattern-match on the verdict and print the schema above.
    The trichotomy label distinguishes $\mathrm{LAMAN}$ from
    $\mathrm{SPARSE\_NOT\_TIGHT}$ on the $\mathrm{.accept}$
    constructor via the constant-time edge-count equation
    $|E(G)| + 3 = 2|V|$.
\end{enumerate}

No blueprint nodes for the binary itself --- it is not Lean-formal
content but the natural delivery vehicle for the runtime claim of
\cref{sec:executable-backends}. The
\cref{def:pebbleGameResult} verdict makes the accepting orientation $D$
and the blocking subset $V'$ structurally available at the CLI's
entry point: $D$ from the verdict's $\mathrm{.accept}\,D$
constructor (its $\mathrm{arcs}$ field is the witness data on the
accept branches) and $V'$ from $\mathrm{.reject}\,V'$ (whose proof
fields are unused at $\mathrm{IO}$-time but carry the structural
certificate). Sample
input files live under \texttt{examples/} with the expected output
schema documented as a comment block at the top of each file.

\emph{Special case: planar rigidity (continued).} The CLI binary
at the default parameters $(k, \ell) = (2, 3)$ is, by
\cref{thm:pebble-game-correct} composed with
\cref{thm:rigidityMatroid-indep-iff-rowIndependent}, an
algorithmic decision procedure for membership in the planar
rigidity matroid $\mathcal{R}(V)$
(\cref{def:rigidityMatroid}), specialised to spanning
subgraphs of edge count $2|V| - 3$ to decide the Laman property
(\cref{def:isLaman}). This is the Jacobs--Hendrickson
$(2, 3)$-pebble game~\cite{jacobsHendrickson1997} as a
formally-verified, end-user-executable program.
